% Generated by docs/report/make_process_figure.py -- do not edit by hand.
\documentclass[fontset=none,11pt,linespread=1.0]{ctexart}
\usepackage[paperwidth=33.867cm,paperheight=19.05cm,margin=0cm]{geometry}
\usepackage{tikz,booktabs,array,amssymb}
\usetikzlibrary{arrows.meta,positioning}
\IfFileExists{/mnt/c/Windows/Fonts/msyh.ttc}{%
  \setCJKmainfont{msyh.ttc}[Path=/mnt/c/Windows/Fonts/,BoldFont=msyhbd.ttc,ItalicFont=msyh.ttc]
  \setCJKsansfont{msyh.ttc}[Path=/mnt/c/Windows/Fonts/,BoldFont=msyhbd.ttc,ItalicFont=msyh.ttc]
  \setCJKmonofont{msyh.ttc}[Path=/mnt/c/Windows/Fonts/]
  \setmainfont{msyh.ttc}[Path=/mnt/c/Windows/Fonts/,BoldFont=msyhbd.ttc,ItalicFont=msyh.ttc]
}{}
\IfFileExists{/mnt/c/Windows/Fonts/consola.ttf}{%
  \setmonofont{consola.ttf}[Path=/mnt/c/Windows/Fonts/,BoldFont=consolab.ttf]
}{}
\pagestyle{empty}\parindent=0pt\renewcommand{\arraystretch}{1.18}
\definecolor{cHot}{HTML}{C0392B}
\definecolor{cInk}{HTML}{1F2D3D}
\definecolor{cEx}{HTML}{1F5FA8}
\definecolor{cRead}{HTML}{B8C4D0}
\definecolor{cProbe}{HTML}{E3B878}
\definecolor{cThink}{HTML}{C9A6C8}
\definecolor{cFull}{HTML}{5B9BD5}
\definecolor{cSim}{HTML}{6BB58C}
\begin{document}
\begin{tikzpicture}[x=1cm,y=1cm,
  lyNwell/.style={fill=black!7,draw=black!45,line width=0.3pt,dash pattern=on 1.5pt off 1pt},
  lyNdiff/.style={fill={rgb,255:red,110;green,185;blue,110}},
  lyPdiff/.style={fill={rgb,255:red,222;green,178;blue,105}},
  lyTap/.style={fill={rgb,255:red,95;green,95;blue,70}},
  lyPoly/.style={fill={rgb,255:red,205;green,60;blue,50}},
  lyM1/.style={fill={rgb,255:red,70;green,120;blue,200},fill opacity=0.55},
  lyM2/.style={fill={rgb,255:red,165;green,85;blue,190},fill opacity=0.5},
  lyVia/.style={fill=black},
  box/.style={draw=cInk,line width=0.6pt,rounded corners=2pt,inner sep=4pt,align=center},
  hd/.style={font=\bfseries,text=cInk,anchor=south west,inner sep=0pt},
  arr/.style={-{Stealth[length=6pt,width=5pt]},line width=0.9pt,draw=cInk},
  devlab/.style={font=\scriptsize\ttfamily,inner sep=0pt},
  devw/.style={font=\scriptsize,inner sep=0pt,text=black!70},
  portlab/.style={font=\scriptsize\ttfamily,anchor=west,inner sep=1.5pt},
  tick/.style={font=\scriptsize,anchor=north,inner sep=1pt,text=black!70},
  sbox/.style={draw=cInk,line width=0.6pt,rounded corners=2pt,minimum width=2.75cm,minimum height=1.55cm,
               inner sep=2pt,align=center,font=\small},
  ex/.style={text=cEx,font=\small,align=center,inner sep=1pt},
]
\useasboundingbox (0,0) rectangle (33.867,19.05);

% ------------------------------------------------------------------ title
\node[anchor=west,font=\Large\bfseries,text=cInk,inner sep=0pt] at (0.7,18.45)
  {基准中 agent 如何生成版图、结果如何评分};
\node[anchor=east,font=\small,text=black!65,inner sep=0pt] at (33.2,18.45)
  {实例：八管施密特触发器（SKY130），DeepSeek V4.1 Flash + Claude Code，带辅助流程的一次运行，官方得分 0.734};
\draw[cInk,line width=0.8pt] (0.7,18.05) -- (33.2,18.05);

% ------------------------------------------------------------------ A: how the layout is produced
\node[hd] at (0.7,17.45) {A\quad 画版图的方式：agent 写程序生成版图，再用检查工具的报告修改程序};

\node[hd,font=\small\bfseries] at (0.7,16.55) {① 题目给定};
\node[anchor=north west,inner sep=0pt,font=\footnotesize] at (0.7,16.4) {%
\begin{tabular}{@{}l@{\hspace{5pt}}c@{\hspace{5pt}}c@{\hspace{5pt}}c@{\hspace{7pt}}l@{}}
\textcolor{gray}{器件} & \textcolor{gray}{类型} & \textcolor{gray}{$W/L$} & \textcolor{gray}{$m$} & \textcolor{gray}{D, G, S} \\[1pt]
{M1} & {n} & {1/0.5} & {2} & {m, in, dvss} \\
{\bfseries\color{cHot}M3} & {\bfseries\color{cHot}p} & {\bfseries\color{cHot}1/0.5} & {\bfseries\color{cHot}6} & {\bfseries\color{cHot}m, in, dvdd} \\
{M2} & {n} & {1/0.5} & {1} & {m, out, dvss} \\
{M4} & {p} & {1/0.5} & {3} & {m, out, dvdd} \\
{M5} & {n} & {1/0.5} & {1} & {out, m, dvss} \\
{M6} & {p} & {1/0.5} & {3} & {out, m, dvdd} \\
{Mdum0} & {n} & {1/0.5} & {1} & {m, dvss, dvss} \\
{Mdum1} & {p} & {1/0.5} & {1} & {dvdd, dvdd, m} \\
\end{tabular}
};
\node[anchor=north west,inner sep=0pt,font=\footnotesize,align=left,text=black!75] at (0.7,11.9)
  {另有：电气指标、器件库文档、\\[2pt]DRC / LVS 规则文件、3 小时时限};

\node[hd,font=\small\bfseries] at (7.0,16.55) {② agent 编写的生成脚本（摘录，全文 160 行 Python）};
\node[anchor=north west,inner sep=5pt,fill=black!4,draw=black!25,line width=0.4pt,rounded corners=2pt,
      font=\fontsize{8.6}{11.2}\selectfont\ttfamily,align=left] (code) at (7.0,16.4) {%
DEVICES~=~[\\
~~~~(\textquotedbl{}M1\textquotedbl{},~\textquotedbl{}n\textquotedbl{},~2.0,~\textquotedbl{}m\textquotedbl{},~\textquotedbl{}dvss\textquotedbl{},~\textquotedbl{}in\textquotedbl{}),\\
\textcolor{cHot}{~~~~(\textquotedbl{}M3\textquotedbl{},~\textquotedbl{}p\textquotedbl{},~6.0,~\textquotedbl{}m\textquotedbl{},~\textquotedbl{}dvdd\textquotedbl{},~\textquotedbl{}in\textquotedbl{}),}\\
~~~~...]\\
for~i,~(name,~kind,~w,~left,~right,~gate)~in~enumerate(DEVICES):\\
~~~~ref~=~c.add\_ref(fac(gate\_width=w,~gate\_length=0.5,\\
~~~~~~~~~~~~~~~~~~~~~~~~nf=1,~guard\_ring=False))\\
~~~~...\\
~~~~\_rect(c,~x~-~SD\_RISER\_W~/~2,~y~-~0.17,\\
~~~~~~~~~~x~+~SD\_RISER\_W~/~2,~0.13,~L.met1drawing)\\
...\\
c.add\_label(net,~position=(tx,~up\_y[net]),\\
~~~~~~~~~~~~layer=L.met2label)\\
};
\node[anchor=north west,inner sep=0pt,font=\footnotesize,align=left,text=black!75] at (7.0,10.55)
  {器件：调用器件库的参数化单元；连线：逐个写金属矩形和通孔的坐标；\\[2pt]本例没有用图形界面或自动布线，全部坐标由脚本算出。};

\node[hd,font=\small\bfseries] at (18.55,16.55) {③ 提交的版图\quad{\normalfont\footnotesize 17.0 µm × 10.0 µm；未画接触孔、局部互连、注入层}};
\path[lyNwell] (23.722,11.697) rectangle (24.500,15.477);
\path[lyNwell] (24.651,12.507) rectangle (25.429,14.667);
\path[lyNwell] (25.580,12.507) rectangle (26.357,14.667);
\path[lyNwell] (26.509,13.047) rectangle (27.286,14.127);
\path[lyNwell] (23.722,11.697) rectangle (27.951,15.477);
\path[lyNdiff] (20.103,13.047) rectangle (20.687,14.127);
\path[lyNdiff] (21.032,13.317) rectangle (21.616,13.857);
\path[lyNdiff] (21.961,13.317) rectangle (22.545,13.857);
\path[lyNdiff] (22.890,13.317) rectangle (23.473,13.857);
\path[lyPdiff] (23.819,11.967) rectangle (24.402,15.207);
\path[lyPdiff] (24.748,12.777) rectangle (25.331,14.397);
\path[lyPdiff] (25.677,12.777) rectangle (26.260,14.397);
\path[lyPdiff] (26.606,13.317) rectangle (27.189,13.857);
\path[lyTap] (19.774,13.511) rectangle (19.882,13.663);
\path[lyTap] (27.600,13.511) rectangle (27.708,13.663);
\path[lyPoly] (20.260,12.809) rectangle (20.530,14.365);
\path[lyPoly] (21.189,13.079) rectangle (21.459,14.095);
\path[lyPoly] (22.118,13.079) rectangle (22.388,14.095);
\path[lyPoly] (23.047,13.079) rectangle (23.317,14.095);
\path[lyPoly] (23.976,11.705) rectangle (24.246,15.469);
\path[lyPoly] (24.905,12.515) rectangle (25.175,14.659);
\path[lyPoly] (25.834,12.515) rectangle (26.104,14.659);
\path[lyPoly] (26.762,13.055) rectangle (27.033,14.119);
\path[lyM1] (20.271,14.213) rectangle (20.519,14.338);
\path[lyM1] (20.271,12.836) rectangle (20.519,12.960);
\path[lyM1] (20.120,13.047) rectangle (20.244,14.127);
\path[lyM1] (20.546,13.047) rectangle (20.671,14.127);
\path[lyM1] (21.200,13.943) rectangle (21.448,14.068);
\path[lyM1] (21.200,13.106) rectangle (21.448,13.230);
\path[lyM1] (21.049,13.317) rectangle (21.173,13.857);
\path[lyM1] (21.475,13.317) rectangle (21.599,13.857);
\path[lyM1] (22.129,13.943) rectangle (22.377,14.068);
\path[lyM1] (22.129,13.106) rectangle (22.377,13.230);
\path[lyM1] (21.978,13.317) rectangle (22.102,13.857);
\path[lyM1] (22.404,13.317) rectangle (22.528,13.857);
\path[lyM1] (23.058,13.943) rectangle (23.306,14.068);
\path[lyM1] (23.058,13.106) rectangle (23.306,13.230);
\path[lyM1] (22.906,13.317) rectangle (23.031,13.857);
\path[lyM1] (23.333,13.317) rectangle (23.457,13.857);
\path[lyM1] (23.987,15.318) rectangle (24.235,15.442);
\path[lyM1] (23.987,11.732) rectangle (24.235,11.856);
\path[lyM1] (23.835,11.967) rectangle (23.960,15.207);
\path[lyM1] (24.262,11.967) rectangle (24.386,15.207);
\path[lyM1] (24.915,14.508) rectangle (25.164,14.632);
\path[lyM1] (24.915,12.542) rectangle (25.164,12.666);
\path[lyM1] (24.764,12.777) rectangle (24.888,14.397);
\path[lyM1] (25.191,12.777) rectangle (25.315,14.397);
\path[lyM1] (25.844,14.508) rectangle (26.093,14.632);
\path[lyM1] (25.844,12.542) rectangle (26.093,12.666);
\path[lyM1] (25.693,12.777) rectangle (25.817,14.397);
\path[lyM1] (26.120,12.777) rectangle (26.244,14.397);
\path[lyM1] (26.773,13.968) rectangle (27.022,14.092);
\path[lyM1] (26.773,13.082) rectangle (27.022,13.206);
\path[lyM1] (26.622,13.317) rectangle (26.746,13.857);
\path[lyM1] (27.049,13.317) rectangle (27.173,13.857);
\path[lyM1] (20.017,13.517) rectangle (20.120,13.657);
\path[lyM1] (20.017,11.475) rectangle (20.179,13.657);
\path[lyM1] (20.006,11.475) rectangle (20.190,11.659);
\path[lyM1] (20.671,13.517) rectangle (20.773,13.657);
\path[lyM1] (20.611,11.000) rectangle (20.773,13.657);
\path[lyM1] (20.600,11.000) rectangle (20.784,11.184);
\path[lyM1] (20.325,14.284) rectangle (20.465,15.461);
\path[lyM1] (20.303,15.277) rectangle (20.487,15.461);
\path[lyM1] (20.946,13.517) rectangle (21.049,13.657);
\path[lyM1] (20.946,11.475) rectangle (21.108,13.657);
\path[lyM1] (20.935,11.475) rectangle (21.119,11.659);
\path[lyM1] (21.599,13.517) rectangle (21.702,13.657);
\path[lyM1] (21.540,11.000) rectangle (21.702,13.657);
\path[lyM1] (21.529,11.000) rectangle (21.713,11.184);
\path[lyM1] (21.254,14.014) rectangle (21.394,15.699);
\path[lyM1] (21.232,15.515) rectangle (21.416,15.699);
\path[lyM1] (21.875,13.517) rectangle (21.978,13.657);
\path[lyM1] (21.875,11.475) rectangle (22.037,13.657);
\path[lyM1] (21.864,11.475) rectangle (22.048,11.659);
\path[lyM1] (22.528,13.517) rectangle (22.631,13.657);
\path[lyM1] (22.469,11.000) rectangle (22.631,13.657);
\path[lyM1] (22.458,11.000) rectangle (22.642,11.184);
\path[lyM1] (22.183,14.014) rectangle (22.323,16.411);
\path[lyM1] (22.161,16.228) rectangle (22.345,16.411);
\path[lyM1] (22.804,13.517) rectangle (22.906,13.657);
\path[lyM1] (22.804,11.713) rectangle (22.966,13.657);
\path[lyM1] (22.793,11.713) rectangle (22.977,11.897);
\path[lyM1] (23.457,13.517) rectangle (23.560,13.657);
\path[lyM1] (23.398,11.000) rectangle (23.560,13.657);
\path[lyM1] (23.387,11.000) rectangle (23.571,11.184);
\path[lyM1] (23.112,14.014) rectangle (23.252,15.936);
\path[lyM1] (23.090,15.753) rectangle (23.274,15.936);
\path[lyM1] (23.733,13.517) rectangle (23.835,13.657);
\path[lyM1] (23.733,11.475) rectangle (23.895,13.657);
\path[lyM1] (23.722,11.475) rectangle (23.906,11.659);
\path[lyM1] (24.386,13.517) rectangle (24.489,13.657);
\path[lyM1] (24.327,11.238) rectangle (24.489,13.657);
\path[lyM1] (24.316,11.238) rectangle (24.500,11.421);
\path[lyM1] (24.041,15.388) rectangle (24.181,15.461);
\path[lyM1] (24.019,15.277) rectangle (24.203,15.461);
\path[lyM1] (24.662,13.517) rectangle (24.764,13.657);
\path[lyM1] (24.662,11.475) rectangle (24.824,13.657);
\path[lyM1] (24.651,11.475) rectangle (24.834,11.659);
\path[lyM1] (25.315,13.517) rectangle (25.418,13.657);
\path[lyM1] (25.256,11.238) rectangle (25.418,13.657);
\path[lyM1] (25.245,11.238) rectangle (25.429,11.421);
\path[lyM1] (24.969,14.578) rectangle (25.110,15.699);
\path[lyM1] (24.948,15.515) rectangle (25.131,15.699);
\path[lyM1] (25.591,13.517) rectangle (25.693,13.657);
\path[lyM1] (25.591,11.713) rectangle (25.753,13.657);
\path[lyM1] (25.580,11.713) rectangle (25.763,11.897);
\path[lyM1] (26.244,13.517) rectangle (26.347,13.657);
\path[lyM1] (26.185,11.238) rectangle (26.347,13.657);
\path[lyM1] (26.174,11.238) rectangle (26.357,11.421);
\path[lyM1] (25.898,14.578) rectangle (26.039,15.936);
\path[lyM1] (25.877,15.753) rectangle (26.060,15.936);
\path[lyM1] (26.519,13.517) rectangle (26.622,13.657);
\path[lyM1] (26.519,11.238) rectangle (26.681,13.657);
\path[lyM1] (26.509,11.238) rectangle (26.692,11.421);
\path[lyM1] (27.173,13.517) rectangle (27.276,13.657);
\path[lyM1] (27.114,11.475) rectangle (27.276,13.657);
\path[lyM1] (27.103,11.475) rectangle (27.286,11.659);
\path[lyM1] (26.827,14.038) rectangle (26.968,16.174);
\path[lyM1] (26.806,15.990) rectangle (26.989,16.174);
\path[lyM1] (19.747,11.000) rectangle (19.909,13.695);
\path[lyM1] (19.736,11.000) rectangle (19.920,11.184);
\path[lyM1] (27.573,11.238) rectangle (27.735,13.695);
\path[lyM1] (27.562,11.238) rectangle (27.745,11.421);
\path[lyM2] (19.450,15.288) rectangle (27.597,15.450);
\path[lyM2] (19.450,15.526) rectangle (27.840,15.688);
\path[lyM2] (19.450,15.763) rectangle (28.083,15.925);
\path[lyM2] (19.450,16.001) rectangle (28.326,16.163);
\path[lyM2] (19.450,16.239) rectangle (28.569,16.401);
\path[lyM2] (19.450,11.724) rectangle (27.840,11.886);
\path[lyM2] (19.450,11.486) rectangle (28.083,11.648);
\path[lyM2] (19.450,11.248) rectangle (28.326,11.410);
\path[lyM2] (19.450,11.011) rectangle (28.569,11.173);
\path[lyM2] (27.516,15.288) rectangle (27.678,15.450);
\path[lyM2] (27.759,11.805) rectangle (27.921,15.607);
\path[lyM2] (28.002,11.567) rectangle (28.164,15.844);
\path[lyM2] (28.245,11.329) rectangle (28.407,16.082);
\path[lyM2] (28.488,11.092) rectangle (28.650,16.320);
\path[lyVia] (20.058,11.527) rectangle (20.139,11.608);
\path[lyVia] (20.652,11.051) rectangle (20.733,11.132);
\path[lyVia] (20.355,15.329) rectangle (20.436,15.410);
\path[lyVia] (20.986,11.527) rectangle (21.067,11.608);
\path[lyVia] (21.581,11.051) rectangle (21.662,11.132);
\path[lyVia] (21.284,15.566) rectangle (21.365,15.647);
\path[lyVia] (21.915,11.527) rectangle (21.996,11.608);
\path[lyVia] (22.509,11.051) rectangle (22.590,11.132);
\path[lyVia] (22.212,16.279) rectangle (22.293,16.360);
\path[lyVia] (22.844,11.764) rectangle (22.925,11.845);
\path[lyVia] (23.438,11.051) rectangle (23.519,11.132);
\path[lyVia] (23.141,15.804) rectangle (23.222,15.885);
\path[lyVia] (23.773,11.527) rectangle (23.854,11.608);
\path[lyVia] (24.367,11.289) rectangle (24.448,11.370);
\path[lyVia] (24.070,15.329) rectangle (24.151,15.410);
\path[lyVia] (24.702,11.527) rectangle (24.783,11.608);
\path[lyVia] (25.296,11.289) rectangle (25.377,11.370);
\path[lyVia] (24.999,15.566) rectangle (25.080,15.647);
\path[lyVia] (25.631,11.764) rectangle (25.712,11.845);
\path[lyVia] (26.225,11.289) rectangle (26.306,11.370);
\path[lyVia] (25.928,15.804) rectangle (26.009,15.885);
\path[lyVia] (26.560,11.289) rectangle (26.641,11.370);
\path[lyVia] (27.154,11.527) rectangle (27.235,11.608);
\path[lyVia] (26.857,16.042) rectangle (26.938,16.123);
\path[lyVia] (19.788,11.051) rectangle (19.869,11.132);
\path[lyVia] (27.613,11.289) rectangle (27.694,11.370);
\node[devlab] at (20.395,10.680) {M1};
\node[devw] at (20.395,10.300) {2};
\node[devlab] at (21.324,10.680) {M2};
\node[devw] at (21.324,10.300) {1};
\node[devlab] at (22.253,10.680) {Mdum0};
\node[devw] at (22.253,10.300) {1};
\node[devlab] at (23.182,10.680) {M5};
\node[devw] at (23.182,10.300) {1};
\draw[cHot,line width=1pt] (23.722,11.632) rectangle (24.500,15.542);
\node[devlab,text=cHot] at (24.111,10.680) {M3};
\node[devw] at (24.111,10.300) {6};
\node[devlab] at (25.040,10.680) {M4};
\node[devw] at (25.040,10.300) {3};
\node[devlab] at (25.969,10.680) {M6};
\node[devw] at (25.969,10.300) {3};
\node[devlab] at (26.897,10.680) {Mdum1};
\node[devw] at (26.897,10.300) {1};
\node[devw,anchor=east] at (19.945,10.300) {$W$/µm};
\draw[gray,line width=0.5pt] (19.975,10.000) -- (23.602,10.000);
\node[devw] at (21.788,9.730) {NMOS};
\draw[gray,line width=0.5pt] (23.691,10.000) -- (27.317,10.000);
\node[devw] at (25.504,9.730) {PMOS（n 阱内）};
\draw[gray,line width=0.4pt] (19.450,15.369) -- (19.270,15.224);
\node[portlab,anchor=east] at (19.270,15.224) {in};
\draw[gray,line width=0.4pt] (19.450,15.607) -- (19.270,15.534);
\node[portlab,anchor=east] at (19.270,15.534) {out};
\draw[gray,line width=0.4pt] (19.450,15.844) -- (19.270,15.844);
\node[portlab,anchor=east] at (19.270,15.844) {m};
\draw[gray,line width=0.4pt] (19.450,16.082) -- (19.270,16.154);
\node[portlab,anchor=east] at (19.270,16.154) {dvdd};
\draw[gray,line width=0.4pt] (19.450,16.320) -- (19.270,16.464);
\node[portlab,anchor=east] at (19.270,16.464) {dvss};
\draw[gray,line width=0.4pt] (19.450,11.805) -- (19.270,11.913);
\node[portlab,anchor=east] at (19.270,11.913) {out};
\draw[gray,line width=0.4pt] (19.450,11.567) -- (19.270,11.603);
\node[portlab,anchor=east] at (19.270,11.603) {m};
\draw[gray,line width=0.4pt] (19.450,11.329) -- (19.270,11.293);
\node[portlab,anchor=east] at (19.270,11.293) {dvdd};
\draw[gray,line width=0.4pt] (19.450,11.092) -- (19.270,10.983);
\node[portlab,anchor=east] at (19.270,10.983) {dvss};
\draw[line width=1pt] (27.570,10.750) -- (28.650,10.750);
\node[devw,anchor=north] at (28.110,10.730) {2 µm};
\path[lyNwell] (18.550,9.140) rectangle (18.850,9.380);
\node[devw,anchor=west] at (18.910,9.260) {n 阱};
\path[lyNdiff] (19.880,9.140) rectangle (20.180,9.380);
\node[devw,anchor=west] at (20.240,9.260) {n 有源};
\path[lyPdiff] (21.480,9.140) rectangle (21.780,9.380);
\node[devw,anchor=west] at (21.840,9.260) {p 有源};
\path[lyTap] (23.080,9.140) rectangle (23.380,9.380);
\node[devw,anchor=west] at (23.440,9.260) {体接触};
\path[lyPoly] (24.410,9.140) rectangle (24.710,9.380);
\node[devw,anchor=west] at (24.770,9.260) {多晶硅};
\path[lyM1] (25.740,9.140) rectangle (26.040,9.380);
\node[devw,anchor=west] at (26.100,9.260) {金属1};
\path[lyM2] (27.070,9.140) rectangle (27.370,9.380);
\node[devw,anchor=west] at (27.430,9.260) {金属2};
\path[lyVia] (28.400,9.140) rectangle (28.700,9.380);
\node[devw,anchor=west] at (28.760,9.260) {通孔};
\node[anchor=west,inner sep=0pt,font=\footnotesize,text=cHot] at (18.550,8.820) {网表中 $W$=1 µm、$m$=6 的管子画成 1 个 $W$=6 µm 的单指管，官方 LVS 判为匹配。};

\node[hd,font=\small\bfseries] at (30.35,16.55) {④ 自检};
\node[box,font=\footnotesize,minimum width=2.9cm] (drc) at (31.85,15.75) {KLayout DRC};
\node[box,font=\footnotesize,minimum width=2.9cm] (lvs) at (31.85,14.65) {KLayout LVS};
\node[box,font=\footnotesize,minimum width=2.9cm] (pex) at (31.85,13.55) {Magic 寄生提取};
\node[box,font=\footnotesize,minimum width=2.9cm] (sim) at (31.85,12.45) {ngspice 后仿真};
\draw[arr] (drc) -- (lvs); \draw[arr] (lvs) -- (pex); \draw[arr] (pex) -- (sim);
\node[anchor=north,inner sep=0pt,font=\footnotesize,align=center,text=black!75] at (31.85,11.75)
  {$\circlearrowleft$ 读报告，改脚本，\\[2pt]重新生成\\[6pt]本次运行共 401 次\\[2pt]工具调用};
\draw[arr] (6.35,14.0) -- (code.west |- 0,14.0);
\draw[arr] (code.east |- 0,14.0) -- ++(0.45,0);
\draw[arr] (29.15,14.0) -- (30.2,14.0);

% ------------------------------------------------------------------ timeline of the run
\node[anchor=east,font=\small\bfseries,text=cInk,inner sep=0pt] at (1.95,7.65) {用时};
\path[fill=black!6] (2.200,7.4) rectangle (32.300,7.9);
\path[fill=cRead] (2.200,7.4) rectangle (2.702,7.9);
\path[fill=cProbe] (2.702,7.4) rectangle (8.722,7.9);
\path[fill=cRead] (8.722,7.4) rectangle (12.735,7.9);
\path[fill=cThink] (12.735,7.4) rectangle (17.083,7.9);
\path[fill=cProbe] (17.083,7.4) rectangle (25.109,7.9);
\path[fill=cThink] (25.109,7.4) rectangle (26.614,7.9);
\path[fill=cFull] (26.614,7.4) rectangle (29.123,7.9);
\path[fill=cSim] (29.123,7.4) rectangle (29.792,7.9);
\draw[gray,line width=0.4pt] (2.200,7.4) -- (2.200,7.300000000000001);
\node[tick] at (2.200,7.28) {0};
\draw[gray,line width=0.4pt] (7.217,7.4) -- (7.217,7.300000000000001);
\node[tick] at (7.217,7.28) {30};
\draw[gray,line width=0.4pt] (12.233,7.4) -- (12.233,7.300000000000001);
\node[tick] at (12.233,7.28) {60};
\draw[gray,line width=0.4pt] (17.250,7.4) -- (17.250,7.300000000000001);
\node[tick] at (17.250,7.28) {90};
\draw[gray,line width=0.4pt] (22.267,7.4) -- (22.267,7.300000000000001);
\node[tick] at (22.267,7.28) {120};
\draw[gray,line width=0.4pt] (27.283,7.4) -- (27.283,7.300000000000001);
\node[tick] at (27.283,7.28) {150};
\draw[gray,line width=0.4pt] (32.300,7.4) -- (32.300,7.300000000000001);
\node[tick] at (32.300,7.28) {180};
\node[tick,anchor=north west] at (32.600,7.28) {min};
\draw[cHot,line width=1pt] (27.618,7.3500000000000005) -- (27.618,8.25);
\node[anchor=south east,inner sep=1.5pt,text=cHot,font=\small] at (27.668,7.9) {第 152 分钟（已用去 92\% 的时间）才第一次生成包含全部器件的版图};
\node[anchor=south west,inner sep=1.5pt,font=\small] at (27.718,7.9) {156 LVS 通过，157 提交，165 结束};
\path[fill=cRead] (2.200,6.45) rectangle (2.520,6.720000000000001);
\node[anchor=west,font=\footnotesize,inner sep=1pt] at (2.600,6.585000000000001) {读题面、文档与规则文件};
\path[fill=cProbe] (6.835,6.45) rectangle (7.155,6.720000000000001);
\node[anchor=west,font=\footnotesize,inner sep=1pt] at (7.235,6.585000000000001) {单器件试验（摸清规则如何识别器件）};
\path[fill=cThink] (13.480,6.45) rectangle (13.800,6.720000000000001);
\node[anchor=west,font=\footnotesize,inner sep=1pt] at (13.880,6.585000000000001) {长思考，回复超出长度上限被截断，没有产出操作};
\path[fill=cFull] (21.800,6.45) rectangle (22.120,6.720000000000001);
\node[anchor=west,font=\footnotesize,inner sep=1pt] at (22.200,6.585000000000001) {写整版脚本并按 DRC / LVS 修改};
\path[fill=cSim] (29.450,6.45) rectangle (29.770,6.720000000000001);
\node[anchor=west,font=\footnotesize,inner sep=1pt] at (29.850,6.585000000000001) {寄生提取与后仿真};

\draw[black!30,line width=0.5pt] (0.7,6.15) -- (33.2,6.15);

% ------------------------------------------------------------------ B: scoring
\node[hd] at (0.7,5.7) {B\quad 评分：对提交的版图独立验收，逐级计分};
\node[sbox] (g) at (2.2,3.7) {提交版图\\$+0.10$};
\node[sbox] (l) at (5.6,3.7) {LVS 通过\\$+0.10$};
\node[sbox] (d) at (9.4,4.65) {DRC 通过\\$+0.10$};
\node[sbox,dashed,minimum width=3.3cm,minimum height=1.75cm] (e) at (9.65,2.8)
  {计分资格\\[2pt]\footnotesize 长宽比 $q \le q_{\max}$\\[1pt]\footnotesize 利用率 $\rho \ge 0.1$，非退化形态};
\node[sbox] (a) at (13.75,2.8) {面积\\$0.20\min(1,\rho/\rho^{*})^{2}$};
\node[sbox] (qq) at (16.95,2.8) {工程质量\\$0.10\,Q$};
\node[sbox] (s) at (20.15,2.8) {后仿真\\$0.40\,F$};
\draw[arr] (g) -- (l);
\draw[arr] (l.east) -- ++(0.35,0) |- (d.west);
\draw[arr] (l.east) -- ++(0.35,0) |- (e.west);
\draw[arr] (e) -- (a);
\node[ex] at (2.2,2.65) {0.10};
\node[ex] at (5.6,2.65) {0.10};
\node[ex,anchor=west] at (10.95,4.65) {0.10};
\node[ex] at (9.65,1.6) {$q$=1.70（$q_{\max}$=3）};
\node[ex] at (13.75,1.6) {0.004\\[1pt]\footnotesize $\rho$=0.114，$\rho^{*}$=0.8};
\node[ex] at (16.95,1.6) {0.030\\[1pt]\footnotesize $Q$=0.30};
\node[ex] at (20.15,1.6) {0.40\\[1pt]\footnotesize $F$=1/1};
\node[anchor=west,font=\small,text=cEx,inner sep=0pt] at (0.7,0.7) {蓝色：本例取值，合计 0.734};
\node[anchor=west,font=\footnotesize,inner sep=0pt,align=left,text=black!80] at (12.3,4.7)
  {$\rho$ = 器件有效面积 / 版图外框面积；$Q$：匹配、走线、供电的几何指标；\\[3pt]
   $F$：Magic 提取寄生（本例 56 个 R、35 个 C）后，通过的工况比例；\\[3pt]
   全部通过：LVS、DRC、计分资格和全部工况都通过。};
\node[anchor=west,font=\footnotesize,inner sep=0pt,align=left,text=cHot] at (6.3,0.7)
  {官方后仿真只测 1 个工况（tt，27 °C）；题面要求 9 个工艺角与温度组合，agent 自测通过其中 6 个。};

% ------------------------------------------------------------------ C: auxiliary flow and what is optimised
\draw[black!30,line width=0.5pt] (22.35,0.4) -- (22.35,6.0);
\node[hd] at (22.75,5.7) {C\quad 辅助流程，以及自动优化改的是哪里};
\path[draw=cInk,line width=0.7pt,dashed,rounded corners=3pt,fill=cFull!8] (22.75,2.8) rectangle (33.2,5.45);
\node[anchor=north west,inner sep=0pt,font=\small\bfseries,text=cInk] at (22.95,5.27) {辅助流程（初版人工编写）};
\node[box,fill=white,font=\small,minimum width=2.9cm,minimum height=1.6cm] (ag) at (24.55,3.8)
  {agent\\[2pt]\footnotesize 模型 + Claude Code\\[2pt]\footnotesize 始终不改};
\node[anchor=west,inner sep=0pt,font=\footnotesize,align=left] at (26.35,3.84)
  {·\,草稿可随时送检，只保留通过检查更多的版本\\[3pt]
   ·\,第 140 分钟提醒剩余时间\\[3pt]
   ·\,第 165 分钟停止；仍有检查未过时，\\[1pt]
   \phantom{·\,}另起一个会话修补，最多 15 分钟};
\node[box,font=\small,minimum width=10.45cm,minimum height=1.1cm,anchor=south west] (opt) at (22.75,0.45)
  {自动优化：另一个 agent 读开发集的运行结果，提出修改；\\[2pt]
   \footnotesize 修改后在开发集上与原版对比，更好才采纳};
\draw[arr] (27.0,1.57) -- (27.0,2.8);
\node[anchor=west,inner sep=0pt,font=\footnotesize,text=black!75,align=left] at (27.25,2.18)
  {只改辅助流程：检查、提示、时间安排\\[2pt]本次采纳的两处修改都在“修补”的提示上};
\end{tikzpicture}
\end{document}
